\section{Introduction}

Mathematics tutoring is a demanding setting for language-model systems because useful answers must be correct, pedagogically clear, and adapted to the learner's current state. A model that only emits a final answer is insufficient: students need intermediate reasoning, explanations at the right level of detail, continuity across turns, and a way to avoid small symbolic or arithmetic errors that can invalidate an otherwise useful solution. At the same time, a tutoring demo must be inspectable. When the system uses retrieved examples, a calculator, or stored student context, those components should be visible enough for a developer or evaluator to understand why the tutor answered as it did.

PrismMathQA addresses this setting with a compact agentic architecture. The framework combines a MetaMathQA-based retrieval corpus, a Qwen3-Embedding-0.6B semantic embedder, Chroma vector search, persistent learner memory, an explicit symbolic verification tool, and a Qwen3-4B reasoning model finetuned through self-distillation with AVSD. The design goal is not only to produce an answer, but to expose the right support around the answer: similar worked examples for grounding, learner context for personalization, and a symbolic verifier for computations where exactness matters.

The scope of this report is mathematical reasoning. PrismMathQA uses MetaMathQA for both retrieval corpus construction and AVSD finetuning data, and uses GSM8K, AIME24, AIME25, and HMMT25 for checkpoint reporting.

The central research question for the broader framework is how to make the model behind the tutor better at step-by-step mathematical reasoning. Standard supervised fine-tuning learns from static reasoning traces, but those traces may not match the model's own rollout distribution. Reinforcement learning with verifiable rewards, including GRPO-style training, can reward correct final answers but often supplies sparse outcome-level feedback \citep{shao2024deepseekmath,deepseek2025r1}. Distillation provides dense distribution-level supervision \citep{hinton2015distilling}; on-policy self-distillation is attractive because it trains on the model's own trajectories while using a teacher distribution to provide dense token-level learning signals \citep{agarwal2024onpolicy,zhao2026selfdistilled}. However, a single privileged teacher context, such as a full reference solution, can be too rigid or can leak information that the student will not observe at test time.

We therefore present Adaptive-View Self-Distillation (AVSD) \citep{nguyen2026avsd} as the training method for the PrismMathQA reasoning model. AVSD uses four privileged views of the same training example: a full solution, a partial rationale, a final answer, and a concise hint. It first extracts a stable consensus direction across teacher views and then gates the residual support that is specific to one or a few views. This produces a reconstructed target distribution for reverse-KL on-policy self-distillation. The resulting training signal is dense like distillation, on-policy like RL-style rollout training, and more robust than committing to one privileged context.

This report makes two main contributions. First, it documents the PrismMathQA framework, showing how RAG, learner memory, and symbolic tool calling work together in an agentic math tutor. Second, it presents AVSD as the training method for the reasoning model, deriving the self-distillation objective and analyzing its benchmark behavior on Qwen3-4B reasoning checkpoints.
